\subsection{Identifikasi Masalah dan Kebutuhan Sistem}
\label{subsec:identifikasi-masalah}
Berdasarkan celah fungsional pada arsitektur DecMed yang telah diuraikan pada Subbab \ref{subsec:analisis-sistem-decmed} dan kebutuhan operasional yang tampak pada Subbab \ref{subsubsec:analisis-skenario-layanan-klinis}, Tugas Akhir ini mengidentifikasi tiga kelompok masalah utama yang sekaligus menjadi fokus pembahasan pada Tugas Akhir ini.

\begin{enumerate}[leftmargin=*]
	\item Mekanisme kontrol akses pada DecMed belum mendukung otorisasi yang granular. Hak akses belum dapat menyatakan batasan rinci seperti jenis data, aksi yang diperbolehkan, masa berlaku, maupun ruang lingkup delegasi. Kondisi ini belum selaras dengan kebutuhan layanan kesehatan yang menuntut penerapan prinsip \textit{need-to-know}.
	\item Skema penyimpanan data \textit{off-chain} pada DecMed belum dirancang secara tersegmentasi untuk mendukung penegakan kebijakan akses granular. Ketika akses diberikan kepada suatu objek data terenkripsi, sistem berpotensi memberikan cakupan akses yang lebih luas daripada kebutuhan aktual pengguna.
	\item Pendelegasian hak akses pada DecMed masih bergantung pada keterlibatan langsung pasien sebagai penerbit akses baru. Ketergantungan ini kurang sesuai dengan alur pelayanan kesehatan yang kolaboratif dan dinamis, karena delegasi antar tenaga medis seharusnya dapat berlangsung secara berantai dengan pembatasan yang makin ketat tanpa memerlukan intervensi ulang dari pasien.
\end{enumerate}

Ketiga kelompok masalah tersebut berkorespondensi langsung dengan rumusan masalah pada subbab \ref{sec:rumusan-masalah}, yaitu kebutuhan untuk merancang mekanisme kontrol akses granular, rancangan segmentasi data, dan fitur atenuasi \textit{offline} untuk delegasi akses berantai. Agar keterhubungan antara analisis masalah, solusi, rancangan, implementasi, dan pengujian dapat ditelusuri, masalah tersebut dirinci menggunakan ID \texttt{M-TA} pada Tabel \ref{tab:identifikasi-masalah-decmed}.

\begin{table}[!ht]
	\centering
	\caption{Identifikasi Masalah DecMed}
	\label{tab:identifikasi-masalah-decmed}
	\small
	\begin{tabular}{|p{0.12\textwidth}|p{0.33\textwidth}|p{0.15\textwidth}|p{0.25\textwidth}|}
		\hline
		\textbf{ID}      & \textbf{Masalah}                                                        & \textbf{STRIDE}                          & \textbf{Kebutuhan Terkait}                                                 \\ \hline
		\texttt{M-TA-01} & Token otorisasi DecMed belum granular                                   & \texttt{T01}, \texttt{T05}, \texttt{T08} & \texttt{KF-TA-01}, \texttt{KF-TA-04}--\texttt{KF-TA-09}                    \\ \hline
		\texttt{M-TA-02} & Data RME belum menjadi objek akses yang tersegmentasi                   & \texttt{T03}, \texttt{T05}, \texttt{T06} & \texttt{KF-TA-05}, \texttt{KF-TA-07}, \texttt{KF-TA-08}, \texttt{KF-TA-09} \\ \hline
		\texttt{M-TA-03} & Alur pemberian akses belum mendukung delegasi antaraktor                & \texttt{T04}, \texttt{T08}               & \texttt{KF-TA-01}, \texttt{KF-TA-04}, \texttt{KF-TA-06}                    \\ \hline
		\texttt{M-TA-04} & Delegasi berantai berisiko memperluas hak akses jika tidak dibatasi     & \texttt{T01}, \texttt{T05}, \texttt{T08} & \texttt{KF-TA-04}, \texttt{KF-TA-06}                                       \\ \hline
		\texttt{M-TA-05} & Role \texttt{MedicalPersonnel} terlalu umum untuk akses tulis per aktor & \texttt{T02}, \texttt{T05}, \texttt{T08} & \texttt{KF-TA-05}, \texttt{KF-TA-07}, \texttt{KF-TA-08}, \texttt{KF-TA-09} \\ \hline
		\texttt{M-TA-06} & Pencabutan dan audit delegasi berantai belum cukup eksplisit            & \texttt{T04}, \texttt{T08}               & \texttt{KF-TA-02}, \texttt{KF-TA-03}                                       \\ \hline
	\end{tabular}
\end{table}

Berdasarkan Tabel \ref{tab:identifikasi-masalah-decmed}, penelitian ini tidak hanya memandang keterbatasan DecMed sebagai persoalan token otorisasi, tetapi sebagai rangkaian masalah yang saling bergantung. \texttt{M-TA-01} menunjukkan keterbatasan ekspresi token akses, \texttt{M-TA-02} menunjukkan belum tersedianya objek data yang cukup granular untuk ditegakkan, \texttt{M-TA-03} dan \texttt{M-TA-04} menunjukkan kebutuhan delegasi terbatas, \texttt{M-TA-05} menunjukkan perlunya validasi berdasarkan subrole tenaga medis, sedangkan \texttt{M-TA-06} menunjukkan kebutuhan pencabutan dan audit pada rantai delegasi. ID masalah tersebut digunakan sebagai dasar pemetaan solusi, rancangan, implementasi, dan pengujian pada bagian selanjutnya. Ancaman \texttt{T07} tidak dipetakan lebih lanjut ke kebutuhan sistem karena \textit{availability} PRE dan IPFS berada di luar fokus utama Tugas Akhir ini.

Berdasarkan permasalahan yang telah diidentifikasi di atas, sistem memerlukan kebutuhan fungsional tambahan yang berfokus pada peningkatan granularitas kontrol akses dan dukungan delegasi antaraktor. Kebutuhan ini dirumuskan untuk melengkapi kebutuhan fungsional DecMed sebelumnya, khususnya pada aspek pemberian akses, pembacaan RME, penulisan RME, pencabutan akses, dan pencatatan audit.

Kebutuhan fungsional tambahan pada penelitian ini tidak dimaksudkan untuk menggantikan kebutuhan fungsional DecMed yang telah dirancang sebelumnya. Sebaliknya, kebutuhan tersebut berperan sebagai perluasan agar sistem dapat mendukung skenario layanan kesehatan kolaboratif sebagaimana dijelaskan pada Subsubbab \ref{subsubsec:analisis-skenario-layanan-klinis}. Dalam skenario tersebut, akses awal terhadap RME hanya diberikan langsung oleh pasien kepada petugas rekam medis, kemudian hak akses tersebut dapat didelegasikan secara terbatas kepada aktor tenaga medis yang terlibat dalam proses pelayanan. Aktor tenaga medis dalam Tugas Akhir ini mencakup dokter, perawat, petugas laboratorium, dan apoteker.

Kebutuhan fungsional tambahan dirumuskan dari sudut pandang aktor, sama seperti pola pada kebutuhan fungsional DecMed pada Lampiran \ref{lampiran:kebutuhan-fungsional-decmed}. Perumusan berbasis aktor digunakan untuk memperjelas pihak yang dapat menerima akses, jenis data yang dapat diakses, operasi yang dapat dilakukan, serta batasan akses yang harus ditegakkan. Batasan tersebut mencakup pasien, episode RME, dataset, fungsi klinis, aksi akses, waktu berlaku, penerima delegasi, dan kedalaman delegasi. Rangkuman kebutuhan fungsional tambahan DecMed ditunjukkan pada Tabel \ref{tab:kebutuhan-fungsional-tambahan-decmed}.

\begin{table}[!ht]
	\centering
	\caption{Kebutuhan Fungsional Tambahan DecMed}
	\label{tab:kebutuhan-fungsional-tambahan-decmed}
	\begin{tabular}{|p{2cm}|p{10.5cm}|}
		\hline
		\textbf{ID} & \textbf{Deskripsi}                                                                                                  \\
		\hline
		% (Penyesuaian dari F06 dan F07)
		KF-TA-01    & Pasien dapat memberikan akses awal kepada petugas rekam medis sesuai konteks pelayanan.                             \\ \hline
		% (Penyesuaian dari F09)
		KF-TA-02    & Pasien dapat mencabut akses awal dan akses turunannya yang telah diberikan.                                         \\ \hline
		KF-TA-03    & Pasien dapat melihat riwayat akses dan delegasi terhadap RME miliknya.                                              \\ \hline
		KF-TA-04    & Petugas rekam medis dapat mendelegasikan akses terbatas kepada perawat dan dokter sesuai cakupan akses.             \\ \hline
		KF-TA-05    & Dokter dapat membaca dan menulis segmen klinis sesuai cakupan akses.                                                \\ \hline
		KF-TA-06    & Dokter dapat mendelegasikan akses terbatas kepada tenaga laboratorium dan apoteker sesuai cakupan akses.            \\ \hline
		KF-TA-07    & Perawat dapat membaca dan menulis segmen klinis sesuai cakupan akses.                                               \\ \hline
		KF-TA-08    & Petugas laboratorium dapat membaca dan menulis segmen klinis sesuai cakupan akses terkait pemeriksaan laboratorium. \\ \hline
		KF-TA-09    & Apoteker dapat membaca dan menulis segmen klinis sesuai cakupan akses terkait pelayanan obat.                       \\ \hline
	\end{tabular}
\end{table}

Terdapat penyesuaian pada F06 dan F07 karena dalam skenario layanan kesehatan kolaboratif, akses awal terhadap RME hanya diberikan langsung oleh pasien kepada petugas rekam medis, kemudian hak akses tersebut dapat didelegasikan secara terbatas kepada aktor tenaga medis yang terlibat dalam proses pelayanan, seperti dokter, perawat, petugas laboratorium, dan apoteker. Penyesuaian ini juga berlaku pada F09 karena dalam skenario tersebut, pasien memerlukan kemampuan untuk mencabut akses awal dan akses turunannya yang telah didelegasikan.

Kebutuhan fungsional tambahan tersebut menjadi dasar bagi perancangan solusi pada tahap berikutnya. Secara khusus, kebutuhan ini digunakan untuk merancang integrasi Macaroons sebagai mekanisme token otorisasi, menentukan struktur pembatasan akses melalui \textit{caveats}, menyusun skema segmentasi data RME berbasis unit data klinis, serta merancang alur delegasi berantai antaraktor. Dengan demikian, rancangan sistem tidak hanya mempertahankan fungsi dasar DecMed, tetapi juga menambahkan mekanisme kontrol akses yang lebih granular, terbatas, dan sesuai dengan kebutuhan layanan kesehatan kolaboratif.

% Untuk mengatasi masalah tersebut, kebutuhan fungsional pada \ref{tab:kebutuhan-fungsional-decmed} perlu diperbaharui agar sesuai dengan kebutuhan sistem. Kebutuhan tersebut dirangkum dalam Tabel \ref{tab:kebutuhan-fungsional-tambahan}.

% \begin{table}[!ht]
% 	\centering
% 	\caption{Pemetaan Masalah dengan Kebutuhan Fungsi Tambahan}
% 	\label{tab:kebutuhan-fungsional-tambahan}
% 	\small
% 	\begin{tabular}{|c|p{0.18\textwidth}|p{0.2\textwidth}|p{0.38\textwidth}|}
% 		\hline
% 		\textbf{Masalah} & \textbf{Fungsionalitas}                            & \textbf{Aktor dan Kebutuhan Akses}                                                                                                                                                        & \textbf{Deskripsi Kebutuhan}                                                                                                                                         \\ \hline
% 		1                & Kontrol Akses Granular                             & Petugas administrasi hanya memerlukan akses administratif; dokter memerlukan akses baca/tulis klinis; petugas laboratorium hanya memerlukan akses terbatas sesuai permintaan pemeriksaan. & Sistem harus mampu menerbitkan token otorisasi yang membawa batasan rinci berdasarkan aktor, jenis data, aksi, masa berlaku, dan konteks penggunaan akses.           \\ \hline
% 		2                & Segmentasi Data \textit{Off-Chain} Terenkripsi     & Data administratif, riwayat klinis, dan hasil laboratorium perlu dapat diakses secara selektif oleh aktor yang berbeda.                                                                   & Sistem harus memecah data RME ke dalam unit penyimpanan \textit{off-chain} yang tersegmentasi agar setiap aktor hanya dapat mengakses himpunan data yang relevan.    \\ \hline
% 		3                & Delegasi Berantai dengan atenuasi \textit{offline} & Dokter perlu dapat mendelegasikan sebagian hak akses kepada petugas laboratorium tanpa meminta pasien menerbitkan ulang token.                                                            & Sistem harus memungkinkan tenaga medis mendelegasikan sebagian hak akses secara lokal dengan batasan yang makin ketat tanpa memerlukan intervensi ulang dari pasien. \\ \hline
% 	\end{tabular}
% \end{table}
